% Supplementary figures. Ported from the chapter-review draft; the figure
% PDFs live in figures/ alongside the main-text ones.
\clearpage
\section*{Supplementary Information}
\setcounter{figure}{0}
\renewcommand{\thefigure}{S\arabic{figure}}

\begin{figure}[H]
 \centering
 \includegraphics[width=\linewidth]{figS_view_design.pdf}
 \caption{\textbf{Cross-view consistency depends partly on overlap between the two views.}
 Two view constructions were compared on the same 2,000 profiles using the same pretrained
 encoder: the canonical training construction, in which the two 50\% CpG subsets are drawn
 independently and therefore share approximately half of their sites, and a disjoint
 construction, in which the subsets partition the valid CpGs and share none. \textbf{a}, Mean
 cosine similarity for matched and unmatched pairs. Bars show the mean and error bars the
 standard deviation across five independent random view draws. Note that the vertical axis
 begins at 0.40 so that the near-identical matched values remain distinguishable.
 \textbf{b}, Retrieval accuracy as a function of the number of candidates considered. Lines
 show the mean and shading the minimum--maximum range across the same five draws; the dashed
 line is chance. \textbf{c}, Similarity matrices for one draw under each construction
 (22 representative profiles, shared colour scale). Matched similarity is nearly unchanged
 between constructions (0.984 versus 0.971) whereas top-1 retrieval falls from 73.8\% to
 25.3\%, converging again as more candidates are considered. Because the encoder was trained
 using independently sampled, potentially overlapping views, the disjoint condition is out of
 distribution for it, and these values should be read as a conservative bound on cross-subset
 consistency rather than as a matched comparison.}
 \label{fig:sup-view-design}
\end{figure}

\begin{figure}[H]
 \centering
 \includegraphics[width=\linewidth]{figS_panel_independence.pdf}
 \caption{\textbf{Robustness of pretrained representations to random CpG subsampling.}
 \textbf{a}, Full profiles (grey) and 50\% CpG subsets of the same profiles (blue) projected
 into a shared principal-component space fitted on the full profiles, with a segment joining
 each of 250 randomly chosen profiles to its own subset. \textbf{b}, Per-profile cosine
 similarity between the full-profile embedding and the embedding of a 50\% subset of the same
 profile, against a shuffled-pairing control. \textbf{c}, A ridge age probe fitted on
 full-profile embeddings of training profiles and applied, without any adaptation, to full and
 to 50\%-subset embeddings of 407 held-out profiles; the annotated interval is a paired
 bootstrap 95\% confidence interval on the difference over 5,000 resamples of those profiles.
 Subsets here are random 50\% draws from a single view seed, not the probe sets of a specific
 array generation, so this experiment does not establish independence from a real measurement
 platform. Performance is also not identical: the probe falls from $R^2=0.863$ to
 $R^2=0.835$, a difference whose confidence interval excludes zero. The result therefore
 indicates substantial robustness to which CpGs are observed, not equivalence.}
 \label{fig:sup-panel-independence}
\end{figure}
